\chapter{Referencial teórico}
\label{cap:referencial}

% Conteúdo planejado em ESTRUTURA-TCC1.md. Abertura do capítulo a escrever quando todas as seções estiverem prontas.

\section{Reabilitação neurológica e lesão medular}
\label{sec:ref-reabilitacao}

A lesão medular traumática tem graves consequências para o estado físico, a vida social e a vida profissional das
pessoas afetadas, e sua avaliação neurológica é regida pelo padrão comum ISNCSCI, que foi uniformemente adotado
pelas equipes clínicas e representa o principal indicador de desfecho neurológico \cite[p.~1,~8]{ahuja2017traumatic}.
Entre as consequências secundárias está a osteoporose, que se manifesta de forma rápida e grave após a lesão e
aumenta o risco de fraturas a tal ponto que até metade das pessoas com lesão medular sofre, em algum momento, uma
fratura de baixo impacto \cite[p.~278]{battaglino2012sci}.

A estimulação elétrica funcional (FES) provoca a contração de músculos paralisados ou enfraquecidos por meio de uma
corrente aplicada aos nervos motores periféricos e coordena essas contrações com o objetivo de alcançar uma
determinada função \cite[p.~327]{peckham2005fes}. Como essa técnica depende de neurônios motores inferiores
intactos \cite[p.~329]{peckham2005fes}, nem todos respondem a ela. Em um estudo realizado por um grupo da
Universidade de Brasília, cujo objetivo era definir as condições para a indicação do ciclismo assistido por FES,
57\% dos 14 voluntários com paraplegia responderam à estimulação com uma contração visível
\cite[p.~209,~211]{guimaraes2016fes}. Em uma revisão sistemática que reuniu 92 estudos sobre esse exercício, os
resultados mais frequentemente avaliados foram o estado dos músculos, o desempenho, o condicionamento aeróbico, a
condição óssea, a força e os resultados funcionais e neurológicos \cite[p.~1,~8]{vanderscheer2021fes}, sendo que as
alterações ósseas podem exigir pelo menos um ano de treinamento \cite[p.~10]{vanderscheer2021fes}. A combinação com
realidade virtual, também utilizada no v-FES cycling, foi testada por \citeonline[p.~1]{duffell2019fes} em um
estudo piloto com 11 participantes, no qual o biofeedback com uso de realidade virtual estimulou o esforço
voluntário durante o ciclismo, e o efeito foi medido por meio da escala motora do ISNCSCI antes e depois do
treinamento. Alguns desses resultados possuem um formulário próprio no Bem-Te-Vi: a densitometria registra a
densidade óssea, a ASIA registra o exame neurológico e o eletrodiagnóstico registra a reação muscular à
estimulação.

O acompanhamento, portanto, depende de avaliações repetidas. \citeonline[p.~331]{rauch2008icf} descrevem a
reabilitação como um ciclo iterativo de avaliação, definição de intervenções, intervenção e reavaliação, no qual as
metas são estabelecidas com base na avaliação inicial. Nesse mesmo sentido, \citeonline[p.~9]{ahuja2017traumatic}
recomendam que o primeiro exame segundo o ISNCSCI sirva como valor de referência para comparação ao longo de todo o
acompanhamento. No entanto, a comparação entre duas avaliações, por si só, não indica se a diferença entre elas tem
importância clínica. A diferença mínima clinicamente importante é a menor diferença de pontuação que o paciente
percebe como benéfica e que justificaria uma mudança no tratamento \cite[p.~408]{jaeschke1989mcid}, e, mesmo no
caso do ISNCSCI, ainda não está definido quantos pontos de melhora devem ser considerados clinicamente importantes
\cite[p.~9]{ahuja2017traumatic}. Portanto, embora o aplicativo consiga calcular a diferença entre as avaliações, a
interpretação dessa diferença depende de limiares clínicos e fica a critério da equipe (Seção~\ref{sec:ref-automacao}).

\section{Instrumentos de avaliação}
\label{sec:ref-instrumentos}

Esta seção apresenta, de forma resumida, para cada formulário do Bem-Te-Vi (Tabela~\ref{tab:formularios}), os dados
que o instrumento coleta e os valores que a fonte manda calcular com base neles, pois são justamente essas regras que
a Seção~\ref{sec:automacao} compara com o código do aplicativo. O Quadro~\ref{qua:instrumentos}, no final desta
seção, reúne o conjunto.

\textbf{CIF.} A Classificação Internacional de Funcionalidade, Incapacidade e Saúde organiza a descrição da
funcionalidade em quatro componentes: funções do corpo, estruturas do corpo, atividades e participação e fatores
ambientais \cite{who2001icf}, identificados pelos prefixos b, s, d e e. Os códigos formam uma hierarquia, na qual o
primeiro nível tem um dígito, o segundo nível tem dois e o terceiro e o quarto níveis acrescentam um dígito cada.
Cada código registrado recebe um qualificador que varia de 0 (nenhum problema) a 4 (problema completo), ou pode
assumir os valores 8 (quando as informações são insuficientes para a qualificação) e 9 (quando a categoria não se
aplica). O manual da OMS estabelece que essas informações devem ser agregadas para fins de análise estatística
\cite[p.~13,~26--28]{who2013howtouse}, mas não descreve como obter a nota da categoria com base nas notas das
categorias mais detalhadas. No aplicativo, essa regra se baseia na planilha da equipe (Seção~\ref{sec:cif-oms}).

\textbf{ASIA.} Os padrões internacionais para a classificação neurológica da lesão medular (ISNCSCI), chamados de
ASIA no aplicativo, baseiam-se em dados brutos coletados nos dois lados do corpo: o tato leve e a picada, com notas
de 0 a 2, em 28 dermátomos; a força, com notas de 0 a 5, em dez músculos-chave; e, nos segmentos sacrais, a
contração anal voluntária e a pressão anal profunda. O segmento que não pode ser examinado recebe a marcação NT, de
não testável. Com base nesses dados, o padrão determina, em uma ordem estabelecida, os escores, os níveis sensitivo
e motor de cada lado, o nível neurológico da lesão, o grau na escala de comprometimento da ASIA (AIS) e a zona de
preservação parcial (ZPP). O grau AIS funciona como a árvore de decisão da Figura~\ref{fig:ref-ais}
\cite[p.~9,~12,~15--22]{rupp2021isncsci}.

\begin{figure}[htb]
	\centering
	\caption{Árvore de decisão do grau AIS}
	\label{fig:ref-ais}
	\begin{tikzpicture}[
			font=\footnotesize,
			pergunta/.style={draw, rounded corners=2pt, align=left, text width=8.2cm, inner sep=4pt},
			grau/.style={draw, very thick, circle, minimum size=7mm, inner sep=0pt, font=\footnotesize\bfseries},
			seta/.style={-{Stealth[length=1.8mm]}, semithick},
			rot/.style={font=\scriptsize, inner sep=1.5pt},
			node distance=7mm,
		]
		\node[pergunta] (q1) {O exame é normal em quem tinha déficit anterior?};
		\node[pergunta, below=of q1] (q2) {Há função sensitiva ou motora em S4-5?};
		\node[pergunta, below=of q2] (q3) {Há contração anal voluntária ou, além da sensibilidade em S4-5, alguma
			função motora, de músculo-chave ou não, a mais de três níveis abaixo do nível motor?};
		\node[pergunta, below=of q3] (q4) {Metade ou mais dos músculos-chave abaixo do nível neurológico têm força
			de pelo menos 3?};
		\node[grau, right=10mm of q1] (e) {E};
		\node[grau, right=10mm of q2] (a) {A};
		\node[grau, right=10mm of q3] (b) {B};
		\node[grau, right=10mm of q4] (d) {D};
		\node[grau, below=8mm of q4] (c) {C};
		\draw[seta] (q1) -- node[rot, above] {sim} (e);
		\draw[seta] (q2) -- node[rot, above] {não} (a);
		\draw[seta] (q3) -- node[rot, above] {não} (b);
		\draw[seta] (q4) -- node[rot, above] {sim} (d);
		\draw[seta] (q1) -- node[rot, right] {não} (q2);
		\draw[seta] (q2) -- node[rot, right] {sim} (q3);
		\draw[seta] (q3) -- node[rot, right] {sim} (q4);
		\draw[seta] (q4) -- node[rot, right] {não} (c);
	\end{tikzpicture}
	\fonte{Elaborada pelo autor com base em \citeonline{rupp2021isncsci}.}
\end{figure}
A ZPP indica o segmento mais caudal com alguma função abaixo de um determinado nível e é aplicada apenas quando falta
a função sacral correspondente; caso não haja preservação da função abaixo do nível em questão, registra-se apenas o
próprio nível. Se o NT impedir alguma dessas determinações, o resultado é registrado como ND, não determinável
\cite[p.~3,~18--19]{rupp2021isncsci}. Apesar das regras detalhadas, a classificação incorreta continua sendo um
problema significativo, especialmente em exames com segmentos NT, e \citeonline[p.~453]{schuld2012isncsci}
demonstraram que os algoritmos computacionais que lidam com esses casos alcançam acurácia equivalente à dos
especialistas.

\textbf{GAS.} Na escala de alcance de metas (GAS, do inglês \textit{goal attainment scaling}), cada meta estabelecida
para o paciente é avaliada com pontos de $-2$ a $+2$, sendo que 0 representa o resultado esperado, e pode receber um
peso igual ao produto da importância atribuída pelo paciente pela dificuldade estimada pela equipe. As pontuações das
metas são combinadas no escore T, com média de 50 e desvio-padrão de 10, calculado de acordo com a
Equação~\ref{eq:gas} \cite[p.~449]{kiresuk1968goal}, que se aplica tanto à linha de base quanto ao resultado
alcançado \cite[p.~364--365]{turnerstokes2009gas}:
\begin{equation}
	\label{eq:gas}
	T = 50 + \frac{10 \sum w_i x_i}{\sqrt{(1 - \rho) \sum w_i^2 + \rho \left( \sum w_i \right)^2}}
\end{equation}
\noindent em que $w_i$ é o peso da meta $i$, $x_i$ é a sua pontuação e $\rho$ é a correlação esperada entre as
metas. \citeonline[p.~449]{kiresuk1968goal} consideram suficiente, na maior parte dos casos, assumir
$\rho = 0{,}3$, valor também adotado por \citeonline[p.~364]{turnerstokes2009gas}. Com uma única meta,
$\sum w_i^2$ e $\left( \sum w_i \right)^2$ se igualam, e o valor de $\rho$ deixa de influenciar o resultado.

\textbf{MEEM.} O Mini-Exame do Estado Mental avalia orientação, memória, atenção e linguagem, sendo que a pontuação
total, que chega a no máximo 30 pontos, é a soma dos pontos obtidos em cada item \cite[p.~190]{folstein1975minimental}.
Na versão recomendada para uso no Brasil, \citeonline[p.~777]{brucki2003mem} publicaram a pontuação mediana de 433
indivíduos saudáveis segundo o nível de escolaridade: 20 para analfabetos, 25 para aqueles com 1 a 4 anos de estudo,
26,5 para 5 a 8 anos, 28 para 9 a 11 anos e 29 para mais de 11 anos. Os autores decidiram não estabelecer pontos de
corte, pois estes podem variar dependendo da doença investigada \cite[p.~780]{brucki2003mem}. Esses valores,
portanto, descrevem o desempenho de indivíduos saudáveis, e sua utilização como limites para a classificação do
resultado do paciente é uma decisão de quem aplica o instrumento.

\textbf{Eletrodiagnóstico e densitometria.} No eletrodiagnóstico, o aplicativo registra a reobase, a cronaxia e a
acomodação de cada grupo muscular. A reobase é a corrente limiar para um estímulo de duração infinitamente longa, e a
cronaxia é a duração do pulso que excita o tecido com o dobro da reobase. A mesma fonte descreve a acomodação como a
resistência da célula a um estímulo que aumenta lentamente, o que requer intensidade cada vez maior para causar
excitação \cite[p.~176,~178]{geddes2004chronaxie}; no entanto, não define o índice de acomodação calculado pelo
aplicativo, sobre o qual este trabalho não encontrou nenhuma outra referência. No campo da densitometria, a posição
oficial da International Society for Clinical Densitometry (ISCD) recomenda, em adultos com lesão medular, o uso da
densidade mineral óssea do quadril total, do fêmur distal e da tíbia proximal para diagnosticar osteoporose, prever o
risco de fraturas dos membros inferiores e acompanhar a resposta ao tratamento, desde que existam dados normativos
\cite{morse2019bone}. O índice de massa corporal, calculado no mesmo formulário, é aqui considerado uma definição
matemática, ou seja, a massa dividida pelo quadrado da altura.

\textbf{Anamnese.} Não existe regra de cálculo para a anamnese, e a questão que se coloca é o que estruturar e o que
deixar como texto livre. \citeonline[p.~181]{rosenbloom2011clinical} analisam a tensão entre a expressividade da
narrativa clínica e a documentação estruturada, que facilita a reutilização dos dados, e recomendam que o
profissional escolha o método de acordo com o fluxo de trabalho e o conteúdo do registro, obtendo os dados
reutilizáveis seja a partir da documentação estruturada, seja do processamento posterior do texto.

\begin{quadro}[htb]
	\centering
	\caption{Instrumentos de avaliação: dados coletados e valores derivados previstos pela fonte}
	\label{qua:instrumentos}
	\footnotesize
	\begin{tabularx}{\textwidth}{|>{\raggedright\arraybackslash}p{0.16\textwidth}|>{\raggedright\arraybackslash}X|>{\raggedright\arraybackslash}X|}
		\hline
		Instrumento & Dados coletados & Valores derivados previstos \\ \hline
		CIF & Qualificador de 0 a 4, 8 ou 9 por código & Nenhuma nota agregada por categoria \\ \hline
		ASIA & Sensibilidade (0 a 2) e força (0 a 5) por segmento e lado; função sacral; NT & Escores, níveis, nível
			neurológico, grau AIS e ZPP \\ \hline
		GAS & Importância, dificuldade e pontuação ($-2$ a $+2$) de cada meta & Peso da meta; escores T inicial e
			alcançado \\ \hline
		MEEM & Pontos dos itens & Escore total (0 a 30); medianas por escolaridade como referência \\ \hline
		Eletrodiag\-nóstico & Reobase, cronaxia e acomodação por grupo muscular & Nenhum índice na fonte consultada \\ \hline
		Densitometria & Densidade mineral óssea por região & Diagnóstico de osteoporose, com dados normativos \\ \hline
		Anamnese & Identificação, contato, datas e texto livre & Nenhum \\ \hline
	\end{tabularx}
	\fonte{Elaborado pelo autor a partir das fontes citadas nesta seção.}
\end{quadro}

\section{Engenharia de requisitos}
\label{sec:ref-requisitos}

\subsection{Requisitos, regras de negócio e critérios de aceitação}
\label{sec:ref-req-conceitos}

A norma ISO/IEC/IEEE 29148 apresenta as características de um requisito bem formulado, tais como ser necessário, não
ambíguo, singular, verificável e correto, e recomenda atribuir a cada requisito atributos como identificador único,
prioridade, justificativa e tipo. De acordo com essa norma, um requisito é verificável se for formulado de forma a
permitir comprovar seu cumprimento, e o conjunto de requisitos só é completo se não contiver itens ainda a definir
\cite[p.~12--15]{iso29148:2018}. Os campos que a Seção~\ref{sec:criterios} exige para cada requisito baseiam-se
nesses atributos. No que diz respeito ao tipo, este trabalho segue \citeonline[p.~25]{glinz2007nonfunctional},
segundo o qual um requisito funcional diz respeito ao comportamento do sistema, aos dados e às respostas a estímulos,
enquanto um requisito não funcional é um atributo, como desempenho, segurança ou usabilidade, ou uma restrição, como
aquelas decorrentes da legislação.

\citeonline[p.~167--173]{wiegers2013software} definem como regras de negócio as políticas, leis e normas segundo as
quais a organização opera. Embora não se trate de requisitos de software, elas são uma fonte rica deles, e os autores
as dividem em fatos, restrições, habilitadores de ação, inferências e cálculos, sendo que as regras de cálculo
transformam dados existentes em dados novos por meio de fórmulas ou algoritmos, que muitas vezes são definidos fora
da organização. Neste trabalho, as regras de cálculo são a média das notas dos itens filhos da CIF, definida pela
planilha da equipe, e as pontuações dos demais instrumentos, definidas na literatura citada na
Seção~\ref{sec:ref-instrumentos}, enquanto regras como a que impede a alteração de uma avaliação já concluída são
restrições. Os mesmos autores consideram que um requisito cujo atendimento não possa ser verificado ainda não está
suficientemente claro, que o critério de aceitação deve levar diferentes observadores à mesma conclusão e que os
testes de aceitação devem ser automatizados sempre que possível \cite[p.~347--348]{wiegers2013software}. Nas regras
de cálculo, essas recomendações são atendidas por meio de casos de entrada e de saída esperada que um script executa,
abordagem adotada nos critérios de aceitação da Seção~\ref{sec:especificacao}.

\subsection{Priorização e rastreabilidade}
\label{sec:ref-req-prioridade}

A técnica de priorização mais comum divide os requisitos em grupos, geralmente em três, e funciona melhor quando cada
grupo tem uma definição clara, pois classificações como alta, média e baixa são interpretadas de maneiras diferentes
pelas partes interessadas. Além disso, deixar de atender a um padrão pode resultar em uma penalidade elevada, mesmo para um requisito que tenha pouca importância para o cliente
\cite[p.~73,~76]{berander2005prioritization}. Portanto, os três níveis da Seção~\ref{sec:criterios} são definidos por
condições verificáveis, como a exigência de uma norma ou a correção de um problema de dados, e não pela importância
percebida. A norma ISO/IEC/IEEE 29148 reconhece que a prioridade decorre do consenso entre as partes interessadas
\cite[p.~15]{iso29148:2018}, mas, sem sessões com a equipe, ela é proposta pelo autor e, posteriormente, avaliada
pela equipe no âmbito do TCC 2.

\citeonline[p.~94]{gotel1994traceability} definem a rastreabilidade de requisitos como a capacidade de acompanhar o
ciclo de vida de um requisito em ambas as direções e distinguem a rastreabilidade pré-especificação, que se refere ao
ciclo de vida do requisito antes de sua inclusão na especificação, da rastreabilidade pós-especificação, que se
refere ao que decorre de sua inclusão. Em estudos com mais de 100 profissionais, os autores atribuíram a maioria dos
problemas de rastreabilidade à fase pré-especificação, ou seja, à dificuldade de identificar a origem de cada
requisito. A matriz de rastreabilidade da Seção~\ref{sec:consolidacao} abrange ambas as direções, pois associa cada
requisito à fonte da qual foi derivado e também ao critério de aceitação e ao elemento do modelo em HL7 FHIR que o
atendem.

\subsection{Requisitos obtidos de fontes documentais}
\label{sec:ref-req-fontes}

Sem acesso às partes interessadas, os requisitos podem ser obtidos a partir de artefatos existentes, como o código do
sistema e as normas a ele aplicáveis. A engenharia reversa permite decifrar os requisitos efetivamente implementados
\cite[p.~15]{chikofsky1990reverse}; no entanto, o código por si só não atende às necessidades do analista, e sua
análise só é útil quando orientada por perguntas específicas \cite[p.~241]{sneed1996extracting}. Como o código mostra
o que foi implementado, e não necessariamente o que deveria ter sido, o requisito obtido é comparado a uma fonte de
referência externa, que, neste trabalho, é a regra publicada do instrumento em questão ou a norma
(Seção~\ref{sec:tecnicas}). As normas, por sua vez, descrevem direitos e obrigações em uma terminologia jurídica
complexa e, às vezes, ambígua, que exige esclarecimento e análise antes de se tornar um requisito implementável
\cite[p.~5]{breaux2008analyzing}. Na área da saúde, \citeonline[p.~350--353]{cysneiros2002health} acrescenta que a
elicitação enfrenta limitações que vão desde a disponibilidade das partes interessadas até restrições legais, que um
mesmo termo pode ter significados contraditórios dependendo de quem o utiliza e que os requisitos não funcionais são
numerosos, como a exigência legal de que a documentação médica não contenha rasuras, o que leva à manutenção do
histórico dos dados e ao registro de quem inseriu cada dado, e o controle de acesso de acordo com a função do
profissional de saúde. A falta de sessões com a equipe é uma dessas limitações, e as decisões que dependem do
vocabulário ou de julgamento clínico são registradas como pendências (Seção~\ref{sec:consolidacao}).

\subsection{Qualidade de dados clínicos}
\label{sec:ref-qualidade}

Em uma revisão sobre a qualidade dos dados de prontuários eletrônicos reutilizados em pesquisas,
\citeonline[p.~145]{weiskopf2013dataquality} identificaram cinco dimensões: completude, que se concentra na presença
das informações do paciente no prontuário; correção, se o dado registrado é verdadeiro; concordância, se os dados do
registro são consistentes entre si ou com outra fonte; plausibilidade, se o dado faz sentido em relação ao que mede;
e atualidade, se o dado reflete o estado do paciente naquele momento. \citeonline[p.~5--7]{kahn2016harmonized}
distinguem ainda a verificação, que compara os dados com expectativas internas, da validação, que os compara com
referências externas, e aplicam essa distinção aos valores calculados: na verificação, o IMC calculado pelo banco e
aquele calculado manualmente devem coincidir; já na validação, os resultados dos algoritmos publicados devem
coincidir com valores de referência externos. Neste trabalho, a comparação do cálculo na tela com o cálculo no
servidor é a verificação, e a comparação do cálculo no aplicativo com a regra publicada do instrumento é a validação.

A identificação do paciente merece atenção à parte. Ao analisar 398.939 registros com duplicatas confirmadas,
\citeonline{just2016patientmatching} verificaram que a maior parte das divergências nos campos de nome vinha de erros
de grafia ou da troca de posição entre partes do nome. Essas dimensões e esses estudos servem de base para classificar
os problemas de dados da Seção~\ref{sec:inv-problemas} e as divergências de cálculo da
Seção~\ref{sec:aut-conformidade}.

\section{Interoperabilidade em saúde e HL7 FHIR}
\label{sec:ref-fhir}

\citeonline{lehne2019interoperability} definem a interoperabilidade como a capacidade de dois ou mais sistemas de
trocar informações e utilizar as informações trocadas, descrevendo-a em quatro níveis. O nível técnico garante a
transmissão de dados entre os sistemas. O nível sintático define o formato e a estrutura dos dados. O nível semântico
garante que o significado dos conceitos seja comum, o que é função de terminologias como SNOMED CT e LOINC. O nível
organizacional abrange processos de trabalho, legislação e políticas. Para o SUS, a Portaria GM/MS nº 1.434/2020, no
art. 232, adota definições semelhantes: a interoperabilidade semântica é o uso de modelos de informação e de
vocabulário padronizado, como terminologias e classificações, que garantam a compreensão humana da estrutura da
informação; e a interoperabilidade sintática é o uso de modelos e técnicas computacionais que garantam a troca de
informações padronizadas entre sistemas, sem perda ou alteração de significado \cite{portaria1434:2020}. Essa
portaria distingue o modelo de informação, que é uma representação conceitual da estrutura de informação com
definição semântica de todos os seus elementos, do modelo computacional, que é a representação dessa estrutura em
linguagem de computador. Nesse sentido, a Seção~\ref{sec:modelo-fhir} apresenta um modelo de informação da avaliação
e o modelo computacional correspondente, escrito com os recursos do FHIR.

O HL7 FHIR (\textit{Fast Healthcare Interoperability Resources}) organiza conceitos comuns na área da saúde em
recursos, que, segundo \citeonline{lehne2019interoperability}, são aproximadamente 140 e são acessados e trocados por
meio de tecnologias da web. A versão R4 (4.0.1) é utilizada no âmbito da RNDS e do BR-Core
(\citeauthor{brcore2026}, \citeyear{brcore2026}, \citeyear{rnds2026guia}). A especificação é apresentada como uma
especificação de plataforma, que geralmente precisa ser adaptada a cada contexto de uso, e essas adaptações
determinam quais elementos dos recursos são utilizados, quais são adicionados e qual terminologia se aplica a cada
elemento. Essas adaptações são publicadas em guias de implementação, e cada perfil representa um conjunto de
restrições sobre um recurso. O que não cabe no recurso é incluído como extensão. Cada elemento pode ter extensões, e
a especificação indica que seu uso não deve ser estigmatizado, desde que cada extensão cumpra as regras de definição
e uso \cite{hl7fhirR4}. Os códigos provêm de sistemas de códigos, cada um dos quais é identificado por uma URI e
descrito por um recurso \textit{CodeSystem}.

Os recursos, os elementos e as extensões que este trabalho usa estão no Quadro~\ref{qua:fhir-elementos}. Dois
pontos da especificação orientam o modelo. Cada preenchimento de um formulário, para um paciente e em um
determinado momento, gera um \textit{QuestionnaireResponse} independente, que pode ser validado em relação ao
questionário para confirmar se os itens obrigatórios foram preenchidos e se as respostas respeitam o tipo e a
cardinalidade. A especificação também distingue esse recurso da \textit{Observation}: o primeiro registra a
pergunta feita e a estrutura do formulário, enquanto a segunda lida com o significado do resultado,
independentemente da pergunta que o originou \cite{hl7fhirR4}.

A Figura~\ref{fig:ref-fhir-exemplo} ilustra uma \textit{Observation} de peso corporal extraída dos exemplos da
especificação. O tipo do recurso, a referência ao paciente, a data e o valor com a unidade constituem a parte
sintática. Os códigos do LOINC e do SNOMED CT e a unidade do UCUM constituem a parte semântica. O mesmo dado contém
códigos de mais de um sistema, e o exemplo original inclui também um código local, o que permite utilizar um sistema
próprio sem abandonar um padrão, quando ele existe.

\begin{figure}[htb]
	\centering
	\caption{Exemplo de \textit{Observation} no FHIR R4}
	\label{fig:ref-fhir-exemplo}
	\begin{minipage}{0.82\textwidth}
	\footnotesize
\begin{verbatim}
{
  "resourceType": "Observation",
  "status": "final",
  "code": {
    "coding": [
      { "system": "http://loinc.org", "code": "29463-7",
        "display": "Body Weight" },
      { "system": "http://snomed.info/sct", "code": "27113001",
        "display": "Body weight" }
    ]
  },
  "subject": { "reference": "Patient/example" },
  "effectiveDateTime": "2016-03-28",
  "valueQuantity": { "value": 185, "unit": "lbs",
    "system": "http://unitsofmeasure.org", "code": "[lb_av]" }
}
\end{verbatim}
	\end{minipage}
	\fonte{Adaptada de \citeonline{hl7fhirR4} (exemplo \texttt{observation-example}), sem a narrativa, a
	categoria, o encontro e dois dos quatro códigos.}
\end{figure}

O guia \textit{Structured Data Capture} (SDC) acrescenta ao \textit{Questionnaire} extensões de preenchimento e
cálculo escritas em FHIRPath, que fecham o Quadro~\ref{qua:fhir-elementos}. O guia também só permite que uma
resposta passe ao estado \textit{completed} quando todos os itens obrigatórios habilitados estiverem respondidos
\cite{hl7sdc}. Assim, a regra de cálculo passa a fazer parte da definição do instrumento, e qualquer sistema
compatível com o SDC pode executá-la.

\begin{quadro}[htb]
	\centering
	\caption{Recursos, elementos e extensões do HL7 FHIR usados no modelo}
	\label{qua:fhir-elementos}
	\footnotesize
	\begin{tabularx}{\textwidth}{|>{\raggedright\arraybackslash}p{0.27\textwidth}|>{\raggedright\arraybackslash}X|>{\raggedright\arraybackslash}p{0.3\textwidth}|}
		\hline
		Nome & O que representa & Uso no modelo (Seção~\ref{sec:modelo-fhir}) \\ \hline
		\multicolumn{3}{|l|}{\textbf{Recursos do FHIR R4}} \\ \hline
		\textit{Patient} & Paciente & Paciente avaliado, no perfil do BR-Core \\ \hline
		\textit{Practitioner} & Profissional de saúde & Avaliador, no perfil do BR-Core \\ \hline
		\textit{Questionnaire} & Formulário, com itens agrupados e restrições sobre as respostas & Definição de
			cada instrumento \\ \hline
		\textit{QuestionnaireResponse} & Respostas de um preenchimento do formulário & Cada avaliação \\ \hline
		\textit{Observation} & Resultado, com o seu significado, independente da pergunta que o originou &
			Resultados por capítulo e por área da CIF \\ \hline
		\textit{CodeSystem} & Sistema de códigos, que pode declarar o tipo de hierarquia entre eles & Catálogo da
			CIF \\ \hline
		\multicolumn{3}{|l|}{\textbf{Elementos dos recursos}} \\ \hline
		\textit{subject} (\textit{QuestionnaireResponse}) & Paciente a que as respostas se referem & Ligação da
			avaliação ao paciente \\ \hline
		\textit{author} (\textit{QuestionnaireResponse}) & Quem recebeu e registrou as respostas & Ligação da
			avaliação ao avaliador \\ \hline
		\textit{authored} (\textit{QuestionnaireResponse}) & Data em que as respostas foram registradas & Data de
			registro, separada da data do exame \\ \hline
		\textit{status} (\textit{QuestionnaireResponse}) & Estado da resposta, como \textit{in-progress} ou
			\textit{completed} & Rascunho e avaliação concluída \\ \hline
		\textit{enableWhen} (item do \textit{Questionnaire}) & Exibição condicional de um item, conforme a resposta a
			outro & Itens exclusivos de cada sexo \\ \hline
		\textit{readOnly} (item do \textit{Questionnaire}) & Item que o usuário não edita & Itens calculados e
			sexo \\ \hline
		\textit{repeats} (item do \textit{Questionnaire}) & Item ou grupo que pode se repetir & Blocos repetidos de
			tamanho variável \\ \hline
		\textit{hasMember} (\textit{Observation}) & Observação que agrupa outras, como em um painel de exames & Área
			que agrupa os seus capítulos \\ \hline
		\textit{derivedFrom} (\textit{Observation}) & Medições de que o valor foi obtido & Resultado ligado às
			respostas \\ \hline
		\textit{dataAbsentReason} (\textit{Observation}) & Motivo da falta de um valor & Valor ausente \\ \hline
		\multicolumn{3}{|l|}{\textbf{Extensões do SDC}} \\ \hline
		\textit{calculatedExpression} & Resposta de um item calculada a partir de outras e atualizada durante o
			preenchimento & Médias da CIF e escores \\ \hline
		\textit{initialExpression} e \textit{launchContext} & Valor inicial de um item, inclusive com dados do
			contexto, como o paciente & Sexo preenchido a partir do cadastro \\ \hline
		\textit{ordinalValue} & Peso numérico de cada opção de resposta & Alternativa ao item inteiro, não adotada
			na CIF \\ \hline
	\end{tabularx}
	\fonte{Elaborado pelo autor a partir de \citeonline{hl7fhirR4} e \citeonline{hl7sdc}.}
\end{quadro}

A RNDS, instituída pela mesma portaria (art. 254-A), é a plataforma nacional de interoperabilidade de informações em
saúde entre estabelecimentos públicos e privados e órgãos de gestão. A portaria também considera vantajosos os
padrões abertos, estáveis, de maior adoção e de menor custo e complexidade de implantação (art. 235)
\cite{portaria1434:2020}. A RNDS troca dados por serviços RESTful em FHIR R4, exige o credenciamento do
estabelecimento e publica um guia para cada modelo de informação. Nenhum dos seis modelos publicados trata de
avaliação funcional ou de escalas de reabilitação \cite{rnds2026guia}. Os perfis de uso geral estão no BR-Core, que
adapta o FHIR ao contexto brasileiro e deixa os casos de uso para guias próprios. Duas regras do BR-Core afetam este
trabalho. O perfil de paciente exige CPF ou CNS. Quando falta um dado, o registro usa a extensão
\textit{data-absent-reason} com o motivo, e não códigos de ``sem informação'' no lugar de valores do domínio. O guia
também não tem perfil de \textit{Questionnaire} nem de \textit{QuestionnaireResponse}, então os formulários usam os
recursos-base do R4 \cite{brcore2026}.

A CIF já está se tornando parte dessa tendência em outros países. De acordo com
\citeonline{nieminen2026icf}, o regulamento sobre o Espaço Europeu de Dados de Saúde adota a CIF como padrão para
dados de funcionalidade nos resumos do paciente e nos relatórios de alta, porém muitos países não dispõem de suporte
padronizado para o armazenamento de dados estruturados da CIF no prontuário, e falta seu mapeamento para outras
terminologias e para padrões de informática em saúde. Os autores recomendam adaptar os sistemas para o registro da
CIF por meio de perfis FHIR e citam iniciativas nesse sentido, como o desenvolvimento de perfis na Finlândia e os
modelos baseados em FHIR do projeto PACIO, que utilizam a CIF como referência semântica para a funcionalidade. Na
própria organização HL7, a CIF possui uma URI registrada no HL7 Terminology (\texttt{http://hl7.org/fhir/sid/icf}),
porém a publicação não traz os códigos \cite{hl7tho2025icf}. Portanto, quem modela dados da CIF em FHIR ainda precisa
decidir de onde vêm os códigos e as descrições, questão retomada na Seção~\ref{sec:modelo-fhir}.

\section{Níveis de automação}
\label{sec:ref-automacao}

\citeonline[p.~286]{parasuraman2000automation} sugerem que a automação pode ser aplicada a quatro classes de funções:
a aquisição de informação, a análise de informação, a seleção de decisão e a implementação da ação, sendo que, em
cada classe, o grau de automação varia de baixo a alto. A automação da aquisição significa que a máquina capta e
registra os dados de entrada, enquanto a automação da análise envolve a integração de várias variáveis em um único
valor. A automação da decisão, por sua vez, envolve sistemas com regras condicionais que determinam a escolha quando
certas condições se verificam, com exemplos na medicina. Os autores admitem níveis altos de automação da informação
quando ela é muito confiável e, quando a confiabilidade não pode ser garantida, desde que o operador mantenha acesso
aos dados brutos \cite[p.~288--289,~292]{parasuraman2000automation}.

Em uma metanálise de 18 experimentos baseada nesse modelo, \citeonline[p.~476,~486]{onnasch2014human} constataram que
níveis mais elevados de automação melhoram o desempenho em atividades rotineiras, mas prejudicam o desempenho e a
consciência da situação quando a automação falha, e que esses efeitos negativos se tornam mais prováveis na
transição da análise de informação para a seleção de ação. Por isso, os autores recomendam que o operador seja
envolvido na tomada de decisão. Este trabalho adota essa demarcação na classificação dos valores das escalas
(Seção~\ref{sec:criterios}): o cálculo da pontuação com base nas respostas é uma automação da análise, enquanto a
interpretação dessa pontuação para fins de decisão sobre o paciente cabe ao profissional.

\section{Trabalhos relacionados}
\label{sec:ref-relacionados}

Os trabalhos que mais se aproximam deste abordam três questões: o registro da CIF em sistemas eletrônicos, o cálculo
da classificação da ISNCSCI por algoritmo e os modelos de dados baseados em FHIR. Na revisão sistemática de
\citeonline[p.~964]{maritz2017icf} sobre o uso da CIF em prontuários eletrônicos, das 17 publicações incluídas, 4
abordaram a especificação de requisitos, 7 o desenvolvimento e 6 a implementação. Entre os desafios, os autores
destacaram que a CIF não é estruturada como uma terminologia formal e que seu uso prático exige a redução do número
de códigos. \citeonline{pongpipatpaiboon2020icf} descrevem as etapas para a criação de um sistema de documentação
baseado na CIF na reabilitação após lesão medular. Os autores classificaram nas categorias da CIF os dados de
entrevistas e prontuários de quatro casos e constataram que os instrumentos de avaliação clínica cobriam menos da
metade das categorias dos conjuntos selecionados da CIF, sendo que a menor cobertura ocorreu em atividades e
participação e em fatores ambientais.

No âmbito da ISNCSCI, além do algoritmo de \citeonline{schuld2012isncsci}, que trata dos segmentos NT
(Seção~\ref{sec:ref-instrumentos}), \citeonline[p.~197]{walden2016isncsci} desenvolveram e validaram o algoritmo do
Rick Hansen Institute em cinco fases, com casos de um registro de lesões medulares. Na última fase, as diferenças
entre o algoritmo e a classificação feita manualmente foram causadas por erros dos profissionais, e os autores
concluíram que o algoritmo determina o nível e a gravidade da lesão com base nos dados brutos do exame. Quanto aos
modelos em FHIR, a revisão de escopo de \citeonline{tabari2024fhir} distingue os modelos dinâmicos, baseados em
fluxos de transformação de dados, dos estáticos e cita \textit{Observation}, \textit{Condition} e \textit{Patient}
como os recursos mais utilizados. As limitações mencionadas dizem respeito à integração de dados, à
interoperabilidade, à padronização, ao desempenho e à escalabilidade ou capacidade de generalização.

Este trabalho se aproxima de cada um desses estudos em um ponto. Assim como os trabalhos de especificação revistos
por \citeonline{maritz2017icf}, ele trata dos requisitos para o registro da CIF em um sistema e, da mesma forma que
\citeonline{pongpipatpaiboon2020icf}, baseia-se em fontes documentais, neste caso o código e as planilhas da equipe,
sem as entrevistas. Da ISNCSCI, adota o cálculo da classificação com base nos dados brutos e utiliza o algoritmo do
Praxis para validar a calculadora do aplicativo (Seção~\ref{sec:automacao}). No que diz respeito aos modelos
revisados por \citeonline{tabari2024fhir}, também utiliza a \textit{Observation} para os resultados, mas armazena as
avaliações no \textit{QuestionnaireResponse}, com as regras de cálculo descritas no questionário
(Seção~\ref{sec:modelo-fhir}).
